\documentclass[11pt,a4paper]{article}
\usepackage[utf8]{inputenc}
\usepackage[T1]{fontenc}
\usepackage[french]{babel}
\usepackage{geometry}
\geometry{margin=2.5cm}
\usepackage{booktabs}
\usepackage{listings}
\usepackage{graphicx}

\title{Journal de bord -- Projet 2 partie 2 : le jeu du tee-shirt\\HAI716I}
\author{Neha HAJEE}
\date{}

\begin{document}

\maketitle

\section{Organisation du binôme}
Nous avons choisi le sujet B qui porte sur le jeu du tee-shirt, parmi les trois qui sont proposés. Nous simulons une population d'agents
qui doivent se mettre d'accord sur une convention, la même couleur pour tous, sans avoir de chef et sans vue d'ensemble. Pour cette
partie du projet, nous avons choisi de nous répartir les tâches via des issues GitHub. Pour ma partie, je m'occupe de l'implémentation du moteur
(\texttt{jeu.py}), des expériences et des mesures (\texttt{experiences.py}), des tests du moteur (\texttt{test\_moteur.py}) et des modifications
Voisinage et Gains asymétriques. Zéinab s'occupe de l'implémentation de l'interface (\texttt{interface.py}), des règles de mise à jour
(\texttt{regles.py}), de leurs tests et des modifications Têtus et Mémoire.
En commun, nous avons travaillé sur l'intégration de nos deux parties, le README et le rendu.
Nous avons opté pour cette répartition en particulier car le moteur et les règles ne communiquent que par un contrat commun, l'interface.
C'est donc le même principe que celui que nous avons suivi à la partie 1.


\section{Étape 0 -- Contrat entre moteur et règles}
\subsection{Interface commune}
Le fichier \texttt{interface.py} définit le contrat partagé entre le moteur et les règles. Les constantes \texttt{ROUGE} et
\texttt{BLEU} remplacent les chaînes littérales dans tout le code. La classe \texttt{Observation} encapsule une rencontre du point
de vue d'un agent : sa couleur, la couleur de son partenaire et le gain obtenu. La classe abstraite \texttt{RegleMiseAJour} sert de
classe mère à toutes les règles ; chaque sous-classe redéfinit \texttt{choisir\_couleur(memoire, couleur\_actuelle, rng)}, qui
renvoie la nouvelle couleur. L'implémentation par défaut lève \texttt{NotImplementedError}, et la méthode \texttt{reset} est un
crochet vide destiné aux règles avec état interne. La classe \texttt{ParametresSimulation} regroupe le nombre d'agents, le nombre
maximal de tours et la graine.

La signature de \texttt{choisir\_couleur} ne donne accès qu'à la mémoire locale de l'agent et à sa couleur courante : aucune
information globale sur la population n'est exposée, conformément à la contrainte d'absence de vue d'ensemble. Le générateur
pseudo-aléatoire \texttt{rng} est instancié une seule fois par le moteur et injecté en paramètre, ce qui garantit la
reproductibilité des simulations à graine fixée.

\subsection{Écarts avec le plan initial}
L'issue \#8 prévoyait des couleurs \texttt{"R"}/\texttt{"B"}, une mémoire sous forme de liste de tuples et des règles
implémentées comme fonctions. Le choix final s'est porté sur des classes : l'accès par attribut
(\texttt{obs.couleur\_partenaire}) est plus lisible et moins sujet aux erreurs que l'accès par indice (\texttt{obs[1]}), et
l'héritage depuis une classe mère reprend le modèle de \texttt{Strategie} de la partie 1. La classe \texttt{Agent}, initialement
dupliquée dans l'interface, est désormais définie uniquement dans \texttt{jeu.py} : l'interface se limite au contrat, l'agent
relève du moteur.

\section{Protocole}
\subsection{Paramètres retenus}
\subsection{Graines et reproductibilité}

\section{Étape 1 -- Le jeu de base}
\subsection{Moteur}
\subsubsection{Classe Agent}
La classe \texttt{Agent} modélise un joueur. Ses attributs sont \texttt{couleur} (chaîne \texttt{ROUGE} ou \texttt{BLEU},
le tee-shirt porté), \texttt{memoire} (liste d'objets \texttt{Observation}, initialisée vide et propre à chaque instance) et
\texttt{regle} (instance de \texttt{RegleMiseAJour} injectée par le constructeur). La méthode \texttt{ajout\_observation} ajoute une
observation en fin de mémoire ; \texttt{mettre\_a\_jour(rng)} délègue le choix de la couleur à \texttt{regle.choisir\_couleur} puis
l'affecte. La règle est un attribut d'instance et non une méthode de classe, ce qui permet d'attribuer des règles différentes
selon les agents (nécessaire pour la modification Têtus). La couleur initiale est fournie par l'appelant : la classe n'effectue
aucun tirage aléatoire.

Consigne précisée par l'enseignant : ni règle de décision globale ni mémoire globale. Une même règle peut être utilisée par
tous les agents, puisqu'elle est appliquée à la mémoire locale de chacun. En revanche, une instance de règle partagée ne doit
pas conserver d'état interne, sous peine de constituer une mémoire globale implicite.

\subsubsection{Population}
\subsubsection{Formation des paires}
\subsubsection{Déroulement d'un tour}
\subsubsection{Partie complète et convergence}
\subsection{Tests du moteur}
\subsection{Mesures}
\subsection{Résultats}
\subsection{Quelle couleur l'emporte ?}

\section{Étape 2 -- Nouvelles règles}
\subsection{Intégration des règles}
\subsection{Résultats}
\subsection{Comparaison avec l'étape 1}

\section{Étape 3 -- Modifications du jeu}
\subsection{Choix des modifications}
\subsection{Voisinage}
\subsubsection{Prédiction}
\subsubsection{Mise en œuvre}
\subsubsection{Résultats}
\subsubsection{Analyse de l'écart}
\subsection{Gains asymétriques}
\subsubsection{Prédiction}
\subsubsection{Mise en œuvre}
\subsubsection{Résultats}
\subsubsection{Analyse de l'écart}

\section{Refaire les résultats}

\section{Impasses et changements de cap}
Première version de \texttt{Agent} : couleur stockée dans une liste, paramètre \texttt{self} absent des méthodes, règle
définie comme méthode de classe, constructeur sans paramètres. Ces erreurs ont été corrigées, puis la classe a été adaptée au
contrat final (appel de \texttt{regle.choisir\_couleur(...)}, mémoire composée d'objets \texttt{Observation}). Conclusion : figer
l'interface avant l'implémentation.

Des balises Markdown présentes dans \texttt{interface.py} provoquaient une \texttt{SyntaxError} à l'import. Elles ont été
supprimées. Conclusion : vérifier que chaque module s'importe correctement avant de le pousser sur la branche d'intégration.

\section{Bilan}

\end{document}
